% TOP_PROBLEM: {"problemId": "problem.bott-conjecture-nonnegatively-curved-simply-connected-manifolds-are-rationally-elliptic", "canonicalTitle": "Bott Conjecture: nonnegatively curved simply connected manifolds are rationally elliptic", "releaseRank": 348, "primaryDomain": "geometry_topology", "primaryDomainLabel": "Geometry & topology", "displayStatus": "open", "releaseStatus": "open"}
% !TeX program = lualatex
% ATTEMPT_STATUS: unresolved
\documentclass[11pt]{article}
\usepackage[margin=25mm]{geometry}
\usepackage{fontspec}
\setmainfont{Latin Modern Roman}
\usepackage{amsmath,amssymb,mathtools}
\usepackage{unicode-math}
\setmathfont{Latin Modern Math}
\usepackage{xurl,hyperref}
\hypersetup{colorlinks=true,urlcolor=blue}
\setcounter{secnumdepth}{0}
\setlength{\parindent}{0pt}
\setlength{\parskip}{5pt}
\setlength{\emergencystretch}{3em}
\begin{document}
\section*{\texorpdfstring{Bott Conjecture: nonnegatively curved simply connected manifolds are rationally elliptic}{Bott Conjecture: nonnegatively curved simply connected manifolds are rationally elliptic}}\label{348}
\subsection{Definitions and mathematical statement}
Let $M$ be a compact simply connected manifold without boundary admitting a metric with nonnegative sectional curvature. Let $\Omega M$ be its based loop space and $b_i(\Omega M;F)=\dim_F H_i(\Omega M;F)$. Ziller's ellipticity formulation asks that for every coefficient field $F$ there are constants $C,d$ with
\[
 b_i(\Omega M;F)\le C(i+1)^d\qquad(i\ge0).
\]
The constants may depend on $M,F$. \emph{Terminology qualification:} the source calls this all-fields condition ``elliptic.'' Rational ellipticity usually means $\sum_{j\ge2}\dim_{{\mathbb{Q}}}(\pi_j(M)\otimes{\mathbb{Q}})<\infty$, equivalently the rational loop-homology growth condition for this finite simply connected setting. The catalogue's ``i.e.'' conflates that rational notion with the stronger all-fields wording. The displayed target retains its explicitly stated all-fields demand rather than reducing it silently to $F={\mathbb{Q}}$.
\par\smallskip\textit{Formulation and concept references:} \href{https://arxiv.org/html/math/0701389v3}{[S1]}.

\subsection{Short English statement}
Must the based loop space of every compact simply connected nonnegatively curved manifold have polynomially bounded Betti numbers over every coefficient field?
\subsection{Research attempt: coefficient fields, constructions, and the missing curvature estimate}
\paragraph{1. Scope and source audit.}
This attempt concerns polynomial growth of each sequence $b_i(\Omega M;F)$, for every field $F$, with constants allowed to depend on $M$ and $F$. It does not replace that assertion by rational ellipticity. Ziller's formulation in Section 1 of [S1], checked on 27 September 2026, explicitly uses every field. The examples and quotient constructions in that survey give context, not a theorem proving the conjecture for arbitrary nonnegatively curved manifolds. The arguments below establish elementary closure properties and several familiar examples. The James construction and the homology Serre spectral sequence are used as standard algebraic-topology inputs; none of their conclusions is inferred from a numerical experiment.

\paragraph{2. Products preserve the required kind of bound.}
Suppose connected based spaces $X_1,\ldots,X_r$ have finite-dimensional loop homology in each degree and satisfy
\[
 b_j(\Omega X_a;F)\leq C_a(j+1)^{d_a}\qquad(j\geq0).
\]
The based-loop functor preserves finite products, and the K\"unneth theorem over a field has no Tor term. Consequently
\[
 b_i\left(\Omega\prod_{a=1}^rX_a;F\right)
 =\sum_{j_1+\cdots+j_r=i}\prod_{a=1}^rb_{j_a}(\Omega X_a;F)
 \leq\left(\prod_a C_a\right)
 {i+r-1\choose r-1}(i+1)^{\sum_a d_a}.
\]
For fixed $r$, the binomial factor is bounded by a constant times $(i+1)^{r-1}$. Thus the desired growth condition is closed under finite products. This conclusion holds separately over each field and does not require a common constant across characteristics. It also explains why taking a product of known examples cannot generate a counterexample merely by increasing the number of loop generators: the product imposes commuting relations between factors.

For $m\geq2$, the standard James--Bott--Samelson calculation gives
\[
 H_*(\Omega S^m;F)\cong T_F(x),\qquad |x|=m-1.
\]
Here $T_F(x)$ is the tensor algebra on one generator, with vector-space basis $1,x,x^2,\ldots$. No graded-commutativity assertion is being made. In particular, in odd generator degree one must not impose $x^2=0$. There is exactly one basis element in every degree divisible by $m-1$, and none in the other degrees. Products of spheres therefore satisfy the conjectured bound over every field. For $(S^2)^r$ the exact answer is
\[
 b_i\bigl(\Omega(S^2)^r;F\bigr)={i+r-1\choose r-1}.
\]
These are closed simply connected manifolds carrying the standard product metrics with nonnegative sectional curvature, so this is an actual geometric subfamily of the target.

\paragraph{3. Complex projective spaces without a splitting assumption.}
The Hopf fibration gives a homotopy-fiber sequence
\[
 \Omega S^{2m+1}\longrightarrow\Omega\mathbb{CP}^{m}
 \longrightarrow S^1\qquad(m\geq1).
\]
It is unnecessary to assert a product decomposition of loop spaces. In the homology spectral sequence of this fibration the base has dimension one. For any local coefficient system with fiber vector space $V$, both $H_0(S^1;V)$ and $H_1(S^1;V)$ have dimension at most $\dim_FV$: they are respectively a cokernel and a kernel of the monodromy minus the identity. Thus, even without identifying that monodromy,
\[
 b_i(\Omega\mathbb{CP}^{m};F)
 \leq b_i(\Omega S^{2m+1};F)+b_{i-1}(\Omega S^{2m+1};F)\leq2.
\]
We interpret a Betti number with negative index as zero. The possible spectral-sequence differentials and extensions can only decrease the resulting upper bound on total dimension. This proves the all-fields growth assertion for complex projective spaces, and the product argument proves it for their products with spheres. Their standard metrics are nonnegatively curved.

The use of a homotopy fibration matters. The map to $S^1$ is the connecting map associated with the Hopf bundle; it is not evaluation of a loop at a point. Confusing it with evaluation would produce a different fiber and invalidate the estimate.

\paragraph{4. A fibration criterion with its monodromy hypothesis visible.}
Let $F_0\to E\to B$ be a based fibration of simply connected spaces of finite type, and suppose $\pi_2(B)=0$. Looping gives a homotopy fibration
\[
 \Omega F_0\longrightarrow\Omega E\longrightarrow\Omega B.
\]
Since $\pi_1(\Omega B)=\pi_2(B)=0$, the base of this loop fibration is simply connected. The $E^2$ term of its homology spectral sequence over any field is consequently
\[
 E^2_{p,q}=H_p(\Omega B;F)\otimes_F H_q(\Omega F_0;F).
\]
Each total degree has finitely many terms. Taking dimensions of the associated graded of the abutment gives
\[
 b_i(\Omega E;F)\leq\sum_{p+q=i}b_p(\Omega B;F)b_q(\Omega F_0;F).
\]
If the fiber and base bounds have exponents $d_F,d_B$, this is at most $C_FC_B(i+1)^{d_F+d_B+1}$. This proves a useful extension rule. It is a statement about the topology of a fibration, not a construction of a nonnegatively curved metric on its total space.

For example, every sphere bundle $S^3\to E\to S^4$ has
\[
 b_i(\Omega E;F)\leq
 \#\{(a,b)\in\mathbb Z_{\geq0}^2:2a+3b=i\}\leq i+1.
\]
The estimate is independent of the twisting class and of the field. In those cases where a separate geometric argument supplies nonnegative curvature, the loop-growth conclusion is already verified. It would be circular to assume that every manifold in the conjecture admits such a bundle presentation. Also, if $\pi_2(B)\neq0$, the displayed untwisted $E^2$ term requires an additional argument; the criterion has not silently discarded that issue.

\paragraph{5. Why rational information alone cannot supply the requested result.}
There is a particularly concrete warning outside the class of closed manifolds. Fix a prime $p$ and an integer $k\geq2$, and let $Y$ be the Moore space obtained by attaching a $(k+1)$-cell to $S^k$ by a map of degree $p$. Put $X=\Sigma Y$. Over $\mathbb Q$, its reduced cellular homology vanishes. Over $\mathbb F_p$, $\widetilde H_*(Y;\mathbb F_p)$ has one generator in degree $k$ and one in degree $k+1$, because its cellular boundary becomes zero. The James calculation therefore gives
\[
 H_*(\Omega X;\mathbb Q)=\mathbb Q,
 \qquad
 H_*(\Omega X;\mathbb F_p)=T_{\mathbb F_p}(a,b),
 \quad |a|=k,\quad |b|=k+1.
\]
The latter Hilbert series is
\[
 \sum_{i\geq0}b_i t^i=\frac1{1-t^k-t^{k+1}}.
\]
To prove failure of polynomial growth without an asymptotic theorem, consider all words of length $\ell$. There are $2^\ell$ of them, distributed among only $\ell+1$ degrees between $k\ell$ and $(k+1)\ell$. Some degree in this interval has Betti number at least $2^\ell/(\ell+1)$. A bound $C(i+1)^d$ would imply
\[
 2^\ell\leq C(\ell+1)((k+1)\ell+1)^d
\]
for every $\ell$, which is impossible. Thus polynomial rational loop homology does not imply the all-fields assertion for arbitrary finite simply connected complexes.

This $X$ is neither a closed manifold nor a supplied example with nonnegative curvature. It is not a counterexample to Bott's conjecture. Its role is exact: any proposed proof which passes from characteristic zero to all fields using only finite CW type has a false step. Curvature would have to impose further restrictions excluding this kind of torsion-driven tensor algebra.

\paragraph{6. Where a direct geometric approach stops.}
One might try to use the energy functional on based loops. Morse theory connects the number of geodesic critical points of bounded index to homology of energy sublevels. To obtain the statement in the question, however, one needs a uniform polynomial estimate that survives passage through all energy levels and works with the chosen coefficient field. A bound on the finitely many ordinary Betti numbers of $M$ does not do that. The Moore-space calculation shows how very small cellular data can generate very large loop homology.

Likewise, finite generation of the Pontryagin algebra alone is insufficient: a tensor algebra on two positive-degree generators is finitely generated and has exponential growth. A successful algebraic argument would need relations strong enough to control the growth, and a geometric proof would need to explain why nonnegative sectional curvature supplies them. None of the product or fibration arguments manufactures those relations for an arbitrary $M$.

\paragraph{7. Diagnostics and outcome.}
The accompanying exact diagnostic computes coefficients of $1/(1-t^k-t^{k+1})$ by its integer recurrence, checks the length-word count, and compares product-of-spheres coefficients with the binomial formula in a finite range. These checks detect grading and convolution mistakes; the infinite growth and closure statements follow from the proofs above, not from the finite range.

\textbf{UNRESOLVED.} The attempt proves the all-fields growth condition for the displayed products and fibration subcases, and isolates a genuine logical obstruction to replacing it by a rational statement. The first remaining gap is a curvature-dependent structural or Morse-theoretic bound for an arbitrary closed simply connected nonnegatively curved manifold.

\subsection{Sources}
\begin{itemize}
\item[S1]W. Ziller, Examples of Riemannian Manifolds with non-negative sectional curvature, arXiv:math/0701389v3 [math.DG], 2007.
\url{https://arxiv.org/html/math/0701389v3}.
\textit{Formulation source. Consulted 24 September 2026: Section 1: Hopf curvature conjectures and Bott ellipticity over every coefficient field..}
\end{itemize}
\end{document}
